% Flavor catalog per-process template.
% process_id: CR013
% family: collider_rs
% owner_agent_id: PKA-CR013
% writer_agent_id: null
% checker_agent_id: null
% opus_signoff_id: null
% Canonical metadata sidecar: CR013.yaml
% Status is controlled by the YAML sidecar status_history, not this comment.

\section{\texorpdfstring{High-Mass Diphoton Resonances, \(pp\to X\to\gamma\gamma\)}{High-Mass Diphoton Resonances, pp to X to gamma gamma}}

\subsection*{Process}
This entry covers high-mass diphoton resonance searches,
\(pp\to X\to\gamma\gamma\), with \(X\) interpreted either as a spin-2
Randall--Sundrum graviton benchmark or as a generic spin-0 heavy scalar
benchmark.  The RS-sensitive channel is the spin-2 interpretation, where the
observed diphoton mass spectrum is translated into limits on the first KK
graviton mass \(m_{G^*}\) for fixed values of the dimensionless curvature
coupling \(k/\bar M_{\rm Pl}\).

\subsection*{Current best limit(s)}
The current channel-specific PDG listing for warped-extra-dimension KK
gravitons quotes the CMS full-Run-2 diphoton result as
\[
  m_{G^*} > 4.8~{\rm TeV}
  \quad (95\%~{\rm CL},~k/\bar M_{\rm Pl}=0.1),
\]
for \(pp\to G^*\to\gamma\gamma\)
(\texttt{pdg2025\_extra\_dimensions\_listing.txt}).  The same PDG review
summarizes the CMS 2024 analysis as excluding KK gravitons below
\(2.2\)--\(5.6\) TeV as \(k/\bar M_{\rm Pl}\) is varied from 0.01 to 0.2,
and the ATLAS 2021 full-Run-2 diphoton analysis as giving
\(m_{G^*}>4.5\) TeV at \(k/\bar M_{\rm Pl}=0.1\)
(\texttt{pdg2025\_extra\_dimensions\_review.txt}).  The CMS public-result
page and arXiv text record the same 95\% CL exclusion range for the 2024
diphoton analysis; HEPData provides the corresponding Figure 3 and Figure 4
upper-limit tables for the spin-2 benchmark
(\texttt{cms2024\_exo\_22\_024\_public\_page.html},
\texttt{hepdata2024\_cms\_diphoton\_record.json}).

\subsection*{Relevance to RS / anarchic-flavor pipeline}
This is a direct collider constraint on a neutral spin-2 resonance with
diphoton branching fraction, not a low-energy flavor observable.  In a minimal
RS1-style benchmark, the quoted bound constrains the lightest graviton
excitation \(m_{G^*}\) after fixing \(k/\bar M_{\rm Pl}\), the production
cross section, and the branching fraction to \(\gamma\gamma\).  It does not
directly constrain the KK-gluon mass used in the anarchic-flavor scan unless a
specific RS spectrum map and coupling convention are supplied.

For the present anarchic-flavor pipeline, the reach is therefore diagnostic
rather than leading.  The quark-scan methodology note quotes the live
RS-anarchy median crossing
\(M_{\rm KK}^{\min}(p50,g_* = 3)=47.26~{\rm TeV}\), far above the few-TeV
direct diphoton reach.  A TeV-scale diphoton resonance bound is still useful as
a cross-check on non-anarchic or non-minimal RS scenarios, and it is especially
relevant when gravity or scalar sectors are light while flavor protection or
alignment weakens the low-energy bounds.

\subsection*{Post-2010 developments}
The post-2010 LHC history is a steady increase in energy, luminosity, and
benchmark coverage relative to the original RS graviton phenomenology papers.
CMS \(7\) TeV diphoton data in arXiv:1112.0688 and ATLAS \(7\) TeV diphoton
data in arXiv:1210.8389 established the early LHC resonance limits in the
\(\gamma\gamma\) channel.  ATLAS arXiv:1504.05511 used the \(8\) TeV
diphoton spectrum and raised the \(k/\bar M_{\rm Pl}=0.1\) RS-graviton
limit to \(2.66\) TeV.  CMS arXiv:1809.00327 moved to \(13\) TeV data and
excluded \(m_{G^*}\) below \(2.3\)--\(4.6\) TeV for coupling parameters from
0.01 to 0.2.  ATLAS arXiv:2102.13405 used \(139~{\rm fb}^{-1}\) and excluded
the RS1 benchmark below \(2.2\), \(3.9\), and \(4.5\) TeV for
\(k/\bar M_{\rm Pl}=0.01\), 0.05, and 0.1.  CMS arXiv:2405.09320 is the
current full-Run-2 CMS update, with \(138~{\rm fb}^{-1}\), refined diphoton
background treatments, spin-0 and spin-2 resonance interpretations, and the
strongest reported RS-graviton diphoton limits for \(k/\bar M_{\rm Pl}\ge
0.1\).

\subsection*{Constraint validity / model dependence}
The quoted mass exclusions are benchmark exclusions.  They assume a production
mechanism, a resonance line shape, a width tied to the chosen
\(k/\bar M_{\rm Pl}\), and the branching pattern of the RS graviton model used
by the experiments.  They are not model-independent lower bounds on every RS
KK mass.  In warped models with Standard Model fields propagating in the bulk,
the couplings of the KK graviton to \(ee\), \(\mu\mu\), and \(\gamma\gamma\)
can be suppressed; the PDG review explicitly warns that the dilepton and
diphoton bounds do not apply directly in that case
(\texttt{pdg2025\_extra\_dimensions\_review.txt}).  Spin-0 diphoton limits
are similarly useful only after specifying whether the scalar is a radion,
singlet, heavy Higgs-like state, or another resonance with a definite
production and branching model.

\subsection*{Code coverage in this repo}
\textbf{NO}.  Greps over \texttt{quarkConstraints/}, \texttt{qcd/},
\texttt{flavorConstraints/}, \texttt{neutrinos/}, \texttt{yukawa/},
\texttt{warpConfig/}, \texttt{solvers/}, \texttt{scanParams/}, and
\texttt{tests/} found no diphoton, RS-graviton, LHC-resonance, or collider
direct-search likelihood implementation.  Adjacent evidence is generic only:
\texttt{solvers/bessel.py:5}--\texttt{6} solves KK masses for gauge bosons
and fermions in RS backgrounds, \texttt{scanParams/scan.py:399} and
\texttt{scanParams/scan.py:429} compute and record \(M_{\rm KK}\), and
\texttt{quarkConstraints/deltaf2.py:1}--\texttt{10} plus
\texttt{quarkConstraints/deltaf2.py:452}--\texttt{563} implement the
low-energy \(\Delta F=2\) KK-gluon-inspired exclusion slice.  None of those
paths evaluates \(pp\to G^*\to\gamma\gamma\), cross-section upper limits, or
mass-exclusion curves.

\subsection*{Implementation difficulty}
\textbf{HIGH}.  A live collider constraint would require a model-specific
RS spectrum and branching calculator plus a reinterpretation of LHC upper
limits, or an external collider-recast workflow such as CheckMATE,
MadAnalysis5, Rivet/Contur, or a dedicated simplified-likelihood treatment.
The blocker is not data availability; it is the translation from a 5D point to
production rates, branching ratios, widths, acceptances, and experimental
likelihoods.

\subsection*{Key references}
\texttt{PDG2025\_ExtraDimensionsListing};
\texttt{PDG2025\_ExtraDimensionsReview};
\texttt{CMS2024\_Diphoton\_ArxivFullText};
\texttt{CMS2024\_Diphoton\_PublicPage};
\texttt{CMS2024\_Diphoton\_HEPData};
\texttt{ATLAS2021\_Diphoton\_ArxivFullText};
\texttt{CMS2018\_Diphoton\_ArxivFullText};
\texttt{ATLAS2015\_Diphoton\_ArxivFullText};
\texttt{RandallSundrum1999};
\texttt{DavoudiaslHewettRizzo1999}.
